% Generated by tools/docx2thesis.py from the Word reference list.
\begin{thebibliography}{999}
\bibitem{ref1} IBM Security and Ponemon Institute, Cost of a Data Breach Report 2025. IBM Corp., 2025.
\bibitem{ref2} H. Mishra, V. Chauhan, I. Tyagi, I. Das, and M. Rawat, ``Ransomware Detection in Linux Systems Using eBPF and AI for Dynamic Analysis,'' IPEC J. Sci. Technol., vol. 4, no. 1, Jun. 2025.
\bibitem{ref3} CISA and FBI, ``ESXiArgs Ransomware Virtual Machine Recovery Guidance,'' Advisory AA23-039A, Feb. 2023.
\bibitem{ref4} CISA and FBI, ``\#StopRansomware: Akira Ransomware,'' Advisory AA24-109A, Apr. 2024.
\bibitem{ref5} Fortinet FortiGuard Labs, ``Ransomware Roundup --- Black Basta,'' Jun. 2022.
\bibitem{ref6} SentinelOne Labs, ``IceFire Ransomware Returns: Now Targeting Linux Enterprise Networks,'' Mar. 2023.
\bibitem{ref7} CISA and FBI, ``\#StopRansomware: LockBit 3.0,'' Advisory AA23-075A, Mar. 2023.
\bibitem{ref8} M. Hirano and R. Kobayashi, ``Evasive Ransomware Attacks Using Low-level Behavioral Adversarial Examples,'' arXiv:2508.08656v1, Aug. 2025.
\bibitem{ref9} B. Bringoltz, E. Halperin, R. Feraru, E. Blaichman, and A. Berman, ``CLEAR: Command Level Annotated Dataset for Ransomware Detection,'' NeurIPS 2025.
\bibitem{ref10} A. Sekar, S. G. Kulkarni, and J. Kuri, ``Leveraging eBPF and AI for Ransomware Nose Out,'' arXiv:2406.14020v1, Jun. 2024.
\bibitem{ref11} D. Zhuravchak et al., ``Monitoring Ransomware with Berkeley Packet Filter,'' CPITS-2023-II, Oct. 2023.
\bibitem{ref12} N. Reategui, R. Pletka, and D. Diamantopoulos, ``On the Generalizability of ML- based Ransomware Detection in Block Storage,'' arXiv:2412.21084v1, Dec. 2024.
\bibitem{ref13} A. Brodzik, T. Malec-Kruszy\'{n}ski, W. Niewolski, M. Tkaczyk, K. Bocianiak, and S.-Y. Loui, ``Ransomware Detection Using Machine Learning in the Linux Kernel,'' arXiv:2409.06452v1, Sep. 2024.
\bibitem{ref14} M. E. Ahmed, H. Kim, S. Camtepe, and S. Nepal, ``Peeler: Profiling Kernel-Level Events to Detect Ransomware,'' in Proc. ESORICS 2021, LNCS vol. 12972, pp. 240--260, 2021.
\bibitem{ref15} Y. Kaya, Y. Chen, M. Botacin, S. Saha, F. Pierazzi, L. Cavallaro, D. Wagner, and T. Dumitras, ``ML-Based Behavioral Malware Detection Is Far From a Solved Problem,'' arXiv:2405.06124v2, IEEE SaTML 2025.
\bibitem{ref16} H. J. Hadi, M. Adnan, Y. Cao, F. B. Hussain, N. Ahmad, M. A. Alshara, and Y. Javed, ``iKern: Advanced Intrusion Detection and Prevention at the Kernel Level Using eBPF,'' Technologies, vol. 12, no. 8, art. 122, MDPI, Jul. 2024.
\bibitem{ref17} M. Willers and T. Philippart, ``Ransomware Detection Using Machine Learning with eBPF,'' Offensive Technologies Project Report, 2024.
\bibitem{ref18} S. Kostopoulos, D. Papatsaroucha, I. Kefaloukos, and E. K. Markakis, ``eIDPS: A Real-time eBPF-based and Machine Learning-powered Network Intrusion Detection and Prevention Solution,'' (preprint), 2024--2025.
\bibitem{ref19} R. Kim, J. Ryu, S. Kim, S. Lee, and S. Kim, ``Detecting Cryptojacking Containers Using eBPF-Based Security Runtime and Machine Learning,'' Electronics, vol. 14, no. 6, art. 1208, MDPI, 2025.
\bibitem{ref20} D. Hitaj, G. Pagnotta, F. DeGaspari, L. DeCarli, and L. V. Mancini, ``Minerva: A File-Based Ransomware Detector,'' in Proc. ACM AsiaCCS 2025, 2025.
\bibitem{ref21} M. Cen, F. Jiang, and R. Doss, ``RansoGuard: A RNN-based Framework Leveraging Pre-attack Sensitive APIs for Early Ransomware Detection,'' Computers \& Security, vol. 150, art. 104293, 2025.
\bibitem{ref22} M. Hirano, R. Hodota, and R. Kobayashi, ``RanSAP: An Open Dataset of Ransomware Storage Access Patterns for Training Machine Learning Models,'' Forensic Sci. Int.: Digital Investigation, 2022.
\bibitem{ref23} M. Hirano and R. Kobayashi, ``RanSMAP: Open Dataset of Ransomware Storage and Memory Access Patterns,'' Computers \& Security, 2025.
\bibitem{ref24} C. van Sloun, V. Woeste, K. Wolsing, J. Pennekamp, and K. Wehrle, ``Detecting Ransomware Despite I/\allowbreak{}O Overhead: A Practical Multi-Staged Approach,'' in Proc. 22nd Annu. Network and Distributed System Security Symp. (NDSS '25), San Diego, CA, USA, Feb. 2025, doi: 10.14722/\allowbreak{}ndss.2025.240764.
\bibitem{ref25} D. Gagulic, L. Zumtaugwald, and S. Sahu, ``Ransomware Detection with Machine Learning in Storage Systems,'' Master Project, Dept. of Informatics, Univ. of Zurich, Feb. 2023.
\bibitem{ref26} M. Hirano and R. Kobayashi, ``Machine Learning Based Ransomware Detection Using Storage Access Patterns Obtained From Live-forensic Hypervisor,'' in Proc. IEEE ISNCC, 2019.
\bibitem{ref27} S. R. Davies, R. Macfarlane, and W. J. Buchanan, ``Comparison of Entropy Calculation Methods for Ransomware Encrypted File Identification,'' Entropy, vol. 24, MDPI, 2022.
\bibitem{ref28} S. R. Davies, R. Macfarlane, and W. J. Buchanan, ``NapierOne: A Modern Mixed File Dataset Alternative to Govdocs1,'' Forensic Sci. Int.: Digital Investigation, vol. 40, 2022.
\bibitem{ref29} C.-Y. Yang and R. Sahita, ``Towards a Resilient Machine Learning Classifier --- a Case Study of Ransomware Detection,'' in Proc. Conf. on Applied ML for Information Security, Washington DC, 2019.
\bibitem{ref30} F. Pierazzi, F. Pendlebury, J. Cortellazzi, and L. Cavallaro, "Intriguing Properties of Adversarial ML Attacks in the Problem Space," in Proc. IEEE Symp. Security and Privacy (S\&P), pp. 1332--1349, 2020.
\bibitem{ref31} R. Jordaney, K. Sharad, S. K. Dash, Z. Wang, D. Papini, I. Nouretdinov, and L. Cavallaro, "Transcend: Detecting Concept Drift in Malware Classification Models," in Proc. USENIX Security Symp., pp. 625--642, 2017.
\bibitem{ref32} J. Luo, Z. Zhang, J. Luo, P. Yang, and R. Jing, ``Sequence-based malware detection using a single-bidirectional graph embedding and multi-task learning framework,'' J. Comput. Secur., vol. 31, no. 1, pp. 1--23, IOS Press, 2023, doi: 10.3233/\allowbreak{}JCS-230041.
\bibitem{ref33} P. M. Anand, P. V. S. Charan, H. Chunduri, and S. K. Shukla, ``LARM: Linux Anti Ransomware Monitor,'' Comput. Secur., vol. 159, 2025, doi: 10.1016/\allowbreak{}j.cose.2025.104700.
\bibitem{ref34} Y. Wang, J. Zhang, H. Huang, X. Wang, X. Wu, S. Pei, L. Zhang, J. Hu, J. Cao, P. Zhang, X. Liao, and H. Jin, ``Ransom Access Memories: Achieving Practical Ransomware Protection in Cloud with DeftPunk,'' in Proc. 18th USENIX Symp. Operating Systems Design and Implementation (OSDI '24), 2024.
\bibitem{ref35} C. Zhu, Z. Hu, K. Mu, R. Sun, L. Han, S. Liu, and J. Yan, ``Minding the Semantic Gap for Effective Storage-Based Ransomware Defense,'' in Proc. 40th IEEE Int. Conf. Massive Storage Systems and Technology (MSST), 2024.
\bibitem{ref36} D. Ma, S. Wang, Y. Zhang, K. Yang, J. Wang, J. Zhu, and Y. Zhang, ``Traveling the Hypervisor and SSD: A Tag-Based Approach Against Crypto Ransomware with Fine-Grained Data Recovery,'' in Proc. 30th ACM SIGSAC Conf. Computer and Communications Security (CCS), 2023.
\bibitem{ref37} K. Higuchi and R. Kobayashi, ``Real-Time Defense System Using eBPF for Machine Learning-Based Ransomware Detection Method,'' in Proc. 2023 Int. Conf. Computational Science and Its Applications Workshops (CANDARW), IEEE, 2023.
\bibitem{ref38} A. Kharraz and E. Kirda, ``Redemption: Real-Time Protection Against Ransomware at End-Hosts,'' in Research in Attacks, Intrusions, and Defenses (RAID), Springer LNCS, 2017.
\bibitem{ref39} A. Continella, A. Guagnelli, G. Zingaro, G. De Pasquale, A. Barenghi, S. Maggi, ``ShieldFS: A Self-Healing, Ransomware-Aware Filesystem,'' in Proc. 32nd Annu. Computer Security Applications Conf. (ACSAC), ACM, 2016, pp. 336--347.
\bibitem{ref40} The Linux man-pages project, ``Linux Programmer's Manual: inotify(7), fanotify(7), auditd(8), strace(1),'' Accessed: Jul. 13, 2026. [Online]. Available: https:/\allowbreak{}/\allowbreak{}man7.org/\allowbreak{}linux/\allowbreak{}man-pages/\allowbreak{}
\bibitem{ref41} B. Gregg, Systems Performance: Enterprise and the Cloud, 2nd ed. Boston, MA, USA: Addison-Wesley Professional, 2020.
\bibitem{ref42} Z. J. Wang, E. Montoya, D. Munechika, H. Yang, B. Hoover, and D. H. Chau, ``DiffusionDB: A Large-scale Prompt Gallery Dataset for Text-to-Image Generative Models,'' arXiv:2210.14896, 2022.
\bibitem{ref43} S. Hochreiter and J. Schmidhuber, ``Long Short-Term Memory,'' Neural Computation, vol. 9, no. 8, pp. 1735--1780, 1997.
\bibitem{ref44} N. Srivastava, G. Hinton, A. Krizhevsky, I. Sutskever, and R. Salakhutdinov, ``Dropout: A Simple Way to Prevent Neural Networks from Overfitting,'' J. Machine Learning Research, vol. 15, pp. 1929--1958, 2014.
\bibitem{ref45} R. J. Williams and J. Peng, ``An Efficient Gradient-Based Algorithm for On-line Training of Recurrent Network Trajectories,'' Neural Computation, vol. 2, no. 4, pp. 490--501, 1990.
\bibitem{ref46} F. Pedregosa et al., ``Scikit-learn: Machine Learning in Python,'' J. Machine Learning Research, vol. 12, pp. 2825--2830, 2011.
\bibitem{ref47} L. Prechelt, ``Early Stopping --- But When?'' in Neural Networks: Tricks of the Trade, LNCS 1524, Springer, 1998.
\bibitem{ref48} T. Fawcett, ``An Introduction to ROC Analysis,'' Pattern Recognition Letters, vol. 27, no. 8, pp. 861--874, 2006.
\bibitem{ref49} H. He and E. A. Garcia, ``Learning from Imbalanced Data,'' IEEE Trans. Knowledge and Data Engineering, vol. 21, no. 9, pp. 1263--1284, 2009.
\bibitem{ref50} A. Kharraz, W. Robertson, D. Balzarotti, L. Bilge, and E. Kirda, ``Cutting the Gordian Knot: A Look Under the Hood of Ransomware Attacks,'' in Proc. Conf. Detection of Intrusions and Malware, and Vulnerability Assessment (DIMVA), 2015.
\bibitem{ref51} N. Scaife, H. Carter, P. Traynor, and K. R. B. Butler, ``CryptoLock (and Drop It): Stopping Ransomware Attacks on User Data,'' in Proc. 36th IEEE Int. Conf. Distributed Computing Systems (ICDCS), 2016, pp. 303--312, doi: 10.1109/\allowbreak{}ICDCS.2016.46.
\end{thebibliography}
